\documentclass[12pt,letterpaper]{article}



%%
%% Required packages:
%%
\usepackage{etex}
\usepackage{amsmath}
\usepackage{amssymb}
\usepackage{amstext}
\usepackage{amsthm}
\usepackage{fancyhdr,fancybox}
\usepackage{mathrsfs} % need for \mathscr command
\usepackage{multicol}
\usepackage[normalem]{ulem}
%\usepackage{pictex,pictexplus}
% \usepackage{pictexwd}
\usepackage{rotating}
\usepackage{url}
\usepackage{xfrac}
\usepackage{booktabs}
\usepackage{color,soul}
\usepackage{comment}

\usepackage{hyperref}
\hypersetup{
    colorlinks=true,     % false: boxed links; true: colored links
    linkcolor=blue,      % color of internal links
    citecolor=blue,      % color of links to bibliography
    filecolor=blue,      % color of file links
    urlcolor=blue        % color of external links
}


\usepackage{tikz}
\usetikzlibrary{backgrounds,calc}

%%
%% Common formatting commands
%%
\setlength{\parindent}{0in}
\setlength{\textwidth}{7in}
\setlength{\evensidemargin}{-0.25in}
\setlength{\oddsidemargin}{-0.25in}
\setlength{\parskip}{.5\baselineskip}
\setlength{\topmargin}{-0.5in}
\setlength{\textheight}{9in}

%               Problem and Part
\newcounter{problemnumber}
\newcounter{partnumber}
\newcounter{subpartnumber}
\newcommand{\Problem}{%
  \refstepcounter{problemnumber}%
  \setcounter{partnumber}{0}%
  \setcounter{subpartnumber}{0}%
  \item[\makebox{\hfil\textbf{\theproblemnumber.}\hfil}]%
  }
\newcommand{\InProblem}[2][1.6]{%
  \refstepcounter{problemnumber}%
  \setcounter{partnumber}{0}%
  \setcounter{subpartnumber}{0}%
  \textbf{\theproblemnumber.}\ \ \parbox[t]{#1in}{#2}%
}
\newcommand{\InSmallProblem}[1]{\InProblem[1.05]{#1}}
\newcommand{\NoProblem}[1]{%
  \mbox{\hfil\textbf{#1}\hfil}%
  }
\newcommand{\UnProblem}{%
  \refstepcounter{problemnumber}%
  \setcounter{partnumber}{0}%
  \setcounter{subpartnumber}{0}%
  \mbox{\hfil\textbf{\theproblemnumber.}\hfil}%
  }
\newcommand{\InFlexProblem}[1]{%
  \refstepcounter{problemnumber}%
  \setcounter{partnumber}{0}%
  \setcounter{subpartnumber}{0}%
  \textbf{\theproblemnumber.}\ \ #1%
}
%%  Part:
\newcommand{\Part}{%
  \renewcommand{\thepartnumber}{(\alph{partnumber})}%
  \refstepcounter{partnumber}
  \setcounter{subpartnumber}{0}%
  \item[(\alph{partnumber})]%
}
\newcommand{\InPart}[2][1.75]{%
  \renewcommand{\thepartnumber}{(\alph{partnumber})}%
  \refstepcounter{partnumber}%
  \setcounter{subpartnumber}{0}%
  (\alph{partnumber})\ \ \parbox[t]{#1in}{#2}%
}
\newcommand{\UnPart}{%
  \refstepcounter{partnumber}%
  \renewcommand{\thepartnumber}{(\alph{partnumber})}%
  \setcounter{subpartnumber}{0}%
  (\alph{partnumber})%
}
\newcommand{\InLargePart}[1]{\InPart[3.05]{#1}}
\newcommand{\InFullPart}[1]{\InPart[6.2]{#1}}
\newcommand{\InSmallPart}[1]{\InPart[1.05]{#1}}
%% SubPart:
\newcommand{\SubPart}{%
  \refstepcounter{subpartnumber}%
  \item[(\roman{subpartnumber})]%
  }
\newcommand{\InSubPart}[2][2.25]{%
  \stepcounter{subpartnumber}%
  (\roman{subpartnumber})\ \ \parbox[t]{#1in}{#2}%
}
\newcommand{\InSmallSubPart}[1]{\InSubPart[1.5]{#1}}
\newcommand{\InTinySubPart}[1]{%
  \stepcounter{subpartnumber}%
  (\roman{subpartnumber})\ \ \makebox{#1}%
}
\newcommand{\InMySubPart}[2][2.25]{%
  \stepcounter{subpartnumber}%
  \parbox[t]{#1in}{\makebox[0.3in][l]{(\roman{subpartnumber})}\ \ {#2}}%
}
\newcommand{\TrueFalse}{\textbf{T}\ \ \textbf{F}\ \ }
\newcommand{\TFPart}[1]{% 
  \stepcounter{partnumber}%
  \setcounter{subpartnumber}{0}%
  \item[(\alph{partnumber})] \TrueFalse\parbox[t]{5.5in}{#1}}

\newcommand{\Points}[1]{(#1 points) \ }
\newcommand{\Point}[1]{(#1 point) \ }


%% For Answers:
\newcounter{answernumber}
\newcommand{\Answer}{%
  \refstepcounter{answernumber}%
  \setcounter{partnumber}{0}%
  \setcounter{subpartnumber}{0}%
  \item[\makebox{\hfil\textbf{\theanswernumber.}\hfil}]%
  }


% Setting up the headers for the first page
%\chead{} % will default to what is typed in the file.
\fancyhead[L]{\textsc{Math 22a}}
\fancyhead[R]{\textsc{Fall 2024}}
\fancyfoot[R]{
	\vspace*{\fill}\footnotesize\textcopyright{} \emph{Harvard Math 22a}	
}
\renewcommand{\headrulewidth}{0pt}

%\fancyhead[C]{}
%	\chead{}	
%	\lhead{}
%	\rhead{}
%	\cfoot{}


% Putting copyright on the footer of each page
%\fancypagestyle{secondpage}{
%	\fancyhead[C]{TESTing again} % change this to be blank
%		%\chead{}	
%	\fancyhead[L]{bears} % change this to be blank
%	\fancyhead[R]{dogs} % change this to be blank
%	\fancyfoot[R]{ %keeps the footer the same
%		\vspace*{\fill}\footnotesize\textcopyright{} \emph{Harvard Math 22a}	
%	}
%}
%\renewcommand{\headrulewidth}{0pt}
%\pagestyle{fancy}
\pagestyle{empty}


%\lhead{}
%\rhead{}
%\rfoot{\vspace*{\fill}\footnotesize\textcopyright{} \emph{Harvard Math 22a}}
%\renewcommand{\headrulewidth}{0pt}
%\pagestyle{fancy}




%%
%% Common shortcut commands:
%%
\newcommand{\ds}{\displaystyle}
\newcommand{\sgn}{\operatorname{sgn}}
\renewcommand{\Pr}{\operatorname{P}}
\newcommand{\CPr}[3][{}]{\Pr_{\text{#1}}\left(#2\,|\,#3\right)}

\newcommand{\C}{\mathbb{C}}
\newcommand{\R}{\mathbb{R}}
\newcommand{\Z}{\mathbb{Z}}
\newcommand{\N}{\mathbb{N}}
\newcommand{\Q}{\mathbb{Q}}
\newcommand{\D}{\mathbb{D}}

\newcommand{\rref}{\operatorname{rref}}
\newcommand{\im}{\operatorname{im}}
\newcommand{\rank}{\operatorname{rank}}
\newcommand{\diag}{\operatorname{diag}}
\newcommand{\Span}{\operatorname{span}}
%\renewcommand{\vec}[1]{\mathbf{#1}}
\newcommand{\charpoly}{\mathscr{P}}
\newcommand{\nicefrac}[2]{\sfrac{#1}{#2}} % using xfrac package
\newcommand{\topology}{\mathcal{T}}
\newcommand{\Inn}{\operatorname{Inn}}

\newcommand{\blankbox}[2]{\fbox{\rule{#1}{0in}\rule{0in}{#2}}}

\newenvironment{myindentpar}[1]%
 {\begin{list}{}%
         {\setlength{\leftmargin}{#1}
          \setlength{\rightmargin}{\leftmargin}}%
         \item[]%
 }
 {\end{list}}
\newtheorem{theorem}{Theorem}
\newtheorem*{NStheorem}{Theorem}
\newtheorem{cor}[theorem]{Corollary}
\newtheorem*{claim}{Claim}
\newtheorem{defn}{Definition}

\newenvironment{amatrix}[1]{%
  \left(\begin{array}{@{}*{#1}{c}|c@{}}
}{%
  \end{array}\right)
}

%--------grstep
% For denoting a Gauss' reduction step.
% Use as: \grstep{\rho_1+\rho_3} or \grstep[2\rho_5 \\ 3\rho_6]{\rho_1+\rho_3}
\newcommand{\grstep}[2][\relax]{%
   \ensuremath{\mathrel{
       {\mathop{\longrightarrow}\limits^{#2\mathstrut}_{
                                     \begin{subarray}{l} #1 \end{subarray}}}}}}
\newcommand{\swap}{\leftrightarrow}


\usetikzlibrary{arrows}
\usepackage{wrapfig}
\usepackage{amsfonts,amsmath}

\makeatletter
\renewcommand*\env@matrix[1][*\c@MaxMatrixCols c]{%
	\hskip -\arraycolsep
	\let\@ifnextchar\new@ifnextchar
	\array{#1}}
\makeatother

\def\env@matrix{\hskip -\arraycolsep
	\let\@ifnextchar\new@ifnextchar
	\array{*\c@MaxMatrixCols c}}

\newcommand{\norm}[1]{\left\lVert#1\right\rVert}

\chead{\Large\textbf{Orthogonalization, $\vec\S$6.3, 6.4}}% 

\begin{document}
\thispagestyle{fancy}

Recall from previously:

\begin{itemize}
\item 	A set of vectors $\{ u_1, \dots, u_p\}$ in $\mathbb{R}^n$ is {\bf orthogonal} if $u_k \cdot u_j=0$ for all $j\neq k$.
\item {\bf Theorem 6.3} Let $A$ be an $m\times n$ matrix, then $$(\text{Row}(A))^{\perp} = \text{Null}(A) \text{ and } (\text{Col}(A))^{\perp}= \text{Null}(A^{T}).$$
\item
{\bf Theorem 6.4:} If $S=\{ u_1, \dots, u_p\}$ is an orthogonal set of non-zero vectors in $\mathbb{R}^n$, then $S$ is linearly independent and hence a basis for $\text{Span} \; S$.
\item An {\bf orthogonal basis} for a subspace $W$ of $\mathbb{R}^n$ is a basis for $W$ that is also an orthogonal set.
\item {\bf Theorem 6.5:} Let $\{u_1, \dots,u_p\}$ be an orthogonal basis for a subspace $W$ of $\mathbb{R}^n$. For each $y\in W$, then $$y=c_1 u_1+\dots+c_p u_p$$ where $$c_j=\frac{y\cdot u_j}{u_j \cdot u_j}.$$
\item Let $u$ be a non-zero vector in $\mathbb{R}^n$ and let $v$ be a vector in $\mathbb{R}^n$, then the {\bf orthogonal projection} of $v$ onto $u$, denoted by $\text{proj}_{u}{v}$, is given by $$\text{proj}_{u}{v}=\frac{v\cdot u}{u \cdot u} u$$ Let $L=\text{Span}\{ u\}$, then we also write $\text{proj}_{L}{v}$ for the orthogonal projection of $v$ onto the line defined by the span of $u$; namely $$\text{proj}_{L}{v}=\text{proj}_{u}{v}.$$
\item {\bf Note:} Theorem 6.5 says that $$y=\text{proj}_{u_1}{y}+\text{proj}_{u_2}{y}+\dots+ \text{proj}_{u_p}{y}.$$
\item A set $\{u_1, \dots, u_p\}$ is {\bf orthonormal} if it is orthogonal and $\norm{u_k}=1$ for all $k$.
\item {\bf Theorem 6.6:} An $m\times n$ matrix $U$ has orthonormal columns if and only if $U^{T}U=I_n$.
\item {\bf Theorem 6.7:} Let $U$ be an $m\times n$ matrix with orthonormal columns. Let $x,y \in \mathbb{R}^n$. Then
\begin{enumerate}
\item $\norm{Ux}=\norm{x}$.
\item $(Ux)\cdot (Uy) = x \cdot y$.
\item $(Ux) \cdot (Uy)=0$ if and only if $x \cdot y=0$.
\end{enumerate}
\item An {\bf orthogonal matrix} $U$ is a square matrix where $U^{-1}=U^{T}$. 
\end{itemize}

\newpage

{\bf Orthogonal Decomposition Theorem:} Let $W$ be a subspace of $\mathbb{R}^n$. Then each $y \in \mathbb{R}^n$ can be uniquely written as $$y=\hat{y} + z$$ where $\hat{y} \in W$ and $z \in W^{\perp}$. 

Moreover, if $\{u_1, \dots, u_p\}$ is any orthogonal basis for $W$, then $$\hat{y}= \left(\frac{y \cdot u_1}{u_1 \cdot u_1} \right)u_1+\dots+  \left(\frac{y \cdot u_p}{u_p \cdot u_p} \right)u_p,$$ and $z=y-\hat{y}$.


\begin{defn} The vector $\hat{y}$ is the {\bf orthogonal projection} of $y$ onto the subspace $W$. Written as $$\hat{y}=\text{proj}_{W}{y}.$$
\end{defn} 

\begin{enumerate}

\Problem Prove the Theorem.

\vfill

{\bf Best Approximation Theorem:} Let $W$ be a subspace of $\mathbb{R}^n$. Let $\hat{y}$ be the orthogonal projection of $y$ onto $W$, then $$\norm{y-\hat{y}} < \norm{y-v}$$ for all $v\in W \setminus \{\hat{y}\}$.

\Problem Prove the Theorem.

\vfill

\newpage

{\bf Theorem 6.10:} If $\{ u_1, \dots, u_p\}$ is an orthonormal basis for $W$, then $$\text{proj}_{W}{y} =(y\cdot u_1) u_1 +\dots +(y \cdot u_p) u_p.$$ If $U=[u_1 \; \dots \; u_p]$, then $$\text{proj}_{W}{y} =U U^{T} y.$$ 

\Problem Prove the theorem.

\newpage


{\bf Question:} How can we find an orthogonal basis for a given subspace?


\Problem Let $$W=\left\{\begin{bmatrix}1\\ 1\\ 1\\ 1\\ \end{bmatrix}, \begin{bmatrix}0\\ 1\\ 1\\ 1\\ \end{bmatrix} ,\begin{bmatrix}0\\ 0\\ 1\\ 1\\ \end{bmatrix}\right\}.$$ Find an orthogonal basis for $W$.

\vfill

\newpage

{\bf Gram-Schmidt Orthogonalization.} Given a basis $\{x_1, \dots, x_p\}$ for a non-zero subspace $W$ of $\mathbb{R}^n$, define 
\begin{align*}
v_1 & = x_1\\
v_2 & = x_2 - \left( \frac{x_2 \cdot v_1}{v_1 \cdot v_1} \right) v_1\\
v_3 & = x_3 -\left( \frac{x_3 \cdot v_1}{v_1 \cdot v_1} \right) v_1-\left( \frac{x_3 \cdot v_2}{v_2 \cdot v_2} \right) v_2\\
&\vdots\\
v_p & = x_p -\left( \frac{x_p \cdot v_1}{v_1 \cdot v_1} \right) v_1-\dots -\left( \frac{x_p \cdot v_{p-1}}{v_{p-1} \cdot v_{p-1}} \right) v_{p-1}.
\end{align*}
Then $\{v_1, \dots, v_p\}$ is an orthogonal basis for $W$. In addition $\text{Span}\{v_1,\dots, v_k\}=\text{Span}\{ x_1, \dots, x_k\}$ for all $1 \leq k \leq p$.

\Problem Prove the Theorem.

\vfill

{\bf QR Factorization.} If $A$ is an $m\times n$ matrix with linearly independent columns, then $A=QR$ where $Q$ is an $m\times n$ matrix whose columns form an orthonormal basis for the column space of $A$ and $R$ is an upper triangular $n \times n$ matrix with positive entries along the diagonal. 

\Problem Prove the Theorem.

\vfill

\newpage

\Problem Let $A=\begin{bmatrix} 1 & 0 & 0\\ 1 & 1 & 0\\ 1 & 1 & 1\\ 1 & 1 & 1\\ \end{bmatrix}$. Find QR factorization of $A$.



\end{enumerate}


\begin{comment}
\newpage
\chead{\Large\textbf{Orthogonalization -- Answers / Solutions}}%
\lhead{}
\rhead{}
\thispagestyle{fancy}

\begin{enumerate}
\Answer 
\begin{proof}
Suppose $\{u_1, \dots, u_p\}$ is an orthogonal basis for $W$, then define $$\hat{y}= \left(\frac{y \cdot u_1}{u_1 \cdot u_1} \right)u_1+\dots+  \left(\frac{y \cdot u_p}{u_p \cdot u_p} \right)u_p,$$ and $z=y-\hat{y}$. We first prove that $z$ is orthogonal to $W$. Then $$z \cdot u_k= (y-\hat{y})\cdot u_k= y\cdot u_k - \hat{y} \cdot u_k= y \cdot u_k -\frac{y\cdot u_k}{u_k \cdot u_k} u_k \cdot u_k=.$$ Thus $z$ is orthogonal to a basis for $W$, and hence $z$ is orthogonal to all of $W$.

We need to prove uniqueness. Suppose $y= \hat{y}_1 +z_1$ for $\hat{y}_1 \in W$ and $z_1 \in W^{\perp}$. Thus we have $\hat{y} +z= \hat{y}_1 +z_1$, and hence $\hat{y} -\hat{y}_1= z_1 -z$. Sinice $W$ is a subspace, $\hat{y} -\hat{y}_1$ is in $W$. Since $\hat{y} -\hat{y}_1= z_1 -z$, $\hat{y} -\hat{y}_1$ is in $W^\perp$. Thus $(\hat{y} -\hat{y}_1)\cdot (\hat{y} -\hat{y}_1)=0$, and hence $(\hat{y} -\hat{y}_1)=0$ and $\hat{y}=\hat{y}_1$. Solving for $z$, we get $z=z_1$. Thus $y=\hat{y}+z$ is the unique way to write $y$ as the sum of an element of $W$ and an element of $W^{\perp}$.
\end{proof}
\Answer 
\begin{proof}
	Let $v \in W\setminus \{\hat{y}\}$, then $$\norm{y-v}^2= \norm{y-\hat{y}+\hat{y}-v}^2.$$ Since $\hat{y}-v \in W$ and $y-\hat{y} \in W^{T}$, we get $\hat{y}-v$ and $y-\hat{y}$ are orthogonal. Thus by the Pythagorean theorem, we get $$\norm{y-v}^2= \norm{y-\hat{y}+\hat{y}-v}^2=\norm{y-\hat{y}}^2+\norm{\hat{y}-v}^2.$$ Since $v \in W\setminus \{\hat{y}\}$, we know $\norm{\hat{y}-v}>0$. Thus $$\norm{y-v}^2>\norm{y-\hat{y}}^2,$$ and taking square roots we get $$\norm{y-v}> \norm{y-\hat{y}}.$$
\end{proof}
\Answer 
\begin{proof}
	Since the vectors $u_k$ all have length one, the usual orthogonal projection formula gives $$\text{proj}_{W}{y} =(y\cdot u_1) u_1 +\dots +(y \cdot u_p) u_p.$$
	
	If $U=[u_1\; \dots \; u_n]$, then 
	$$UU^{T}y=\begin{bmatrix} u_1 & \dots & u_p\\ \end{bmatrix} \begin{bmatrix} u_1^{T}\\ \vdots\\ u_p^{T}\\ \end{bmatrix}y=\begin{bmatrix} u_1 & \dots & u_p\\ \end{bmatrix} \begin{bmatrix} u_1\cdot y \\ \vdots\\ u_p\cdot y\\ \end{bmatrix}=(u_1\cdot y)u_1+\dots+(u_p\cdot y)u_p=\text{proj}_{W}{y}.$$
\end{proof}
\Answer This example is from the textbook. The solution presented here is essentially the same. Let $x_1, x_2, x_3$ be the three basis vectors for $W$.

Let $v_1=x_1=\begin{bmatrix} 1\\ 1\\ 1\\ 1\\ \end{bmatrix}$. Now we want a vector $v_2 \in W$ that is orthogonal to $v_1$. Thus we subtract the orthogonal projection of the second basis vector onto $v_1$ to get the orthogonal component of the second basis vector; namely,
$$v_2= x_2 - \text{proj}_{v_1} x_2= \begin{bmatrix} 0\\ 1\\ 1\\ 1\\ \end{bmatrix} -\frac{3}{4}\begin{bmatrix} 1\\ 1\\ 1\\ 1\\ \end{bmatrix} = \begin{bmatrix} -\frac{3}{4}\\[0.5em] \frac{1}{4}\\[0.5em] \frac{1}{4}\\[0.5em] \frac{1}{4}\\ \end{bmatrix}.$$ It seems the subsequent computations will be more pleasant if we scale $v_2$ to have integer entries, so we take $v_2'=3v_2$.

Now we want a third vector $v_3 \in W$ such that $v_2$ is orthogonal to $v_1$ and $v_2$. Again we substract the amount of $x_3$ in the direction of $v_1$ and $v_2$; namely, $$v_3= x_3-\text{proj}_{\text{Span}\{v_1, v_2\}} x_3=x_3-\frac{x_3\cdot v_1}{v_1\cdot v_1} v_1 -\frac{x_3\cdot v_2'}{v_2'\cdot v_2'} v_2'=x_3-\frac{1}{2}v_1-\frac{2}{12}v_2'=\begin{bmatrix}0\\ -\frac{2}{3}\\\frac{1}{3}\\ \frac{1}{3}\\ \end{bmatrix}.$$ Now $v_1, v_2, v_3$ form an orthogonal basis for $W$.
\Answer The following is the proof from Lay.
\begin{proof}
For $1\leq k \leq p$, let $W_k = \text{Span}\{x_1,\dots, x_k\}$. Set $v_1=x_1$. Then $W_1=\text{Span}\{v_1\}$. 

For $k<p$, $\{v_1, \dots, v_k$ is an orthogonal basis for $W_k$. Define $$v_{k+1}= x_{k+1} -\text{proj}_{W_k} x_{k+1}.$$ By the Orthogonal Decomposition Theorem, $v_{k+1}$ is orthogonal to $W_k$ and $v_{k+1}\in W_{k+1}$. Note $v_{k+1}\neq 0$, so $\{ v_1, \dots, v_{k+1}\}$ is an orthogonal basis for $W_{k+1}$. We continue the process until $k+1=p$.
\end{proof}
\Answer 
\begin{proof} The columns of $A$ form a basis for $\text{Col}(A)$. Apply Gram-Schmidt to get an orthonormal basis $u_1, \dots, u_n$ for the column space of $A$. Let $Q=[u_1\; \dots \; u_n]$. For each $k$, $$x_k \in \text{Span}\{u_1, \dots, u_k\},$$ and hence there exists scalars $r_{1k}, \dots r_{kk}$ so that $$x_k = r_{1k}x_1+\dots+ r_{kk} x_k.$$ Now $x_k = Q r_k$ where $r_k =\begin{bmatrix} r_{1k}\\ \vdots\\ r_{kk}\\ 0\\ \vdots\\ 0\\ \end{bmatrix}$. Define $R=[r_1 \; \dots \; r_n]$, then $$A=[x_1 \; \dots \; x_n]=[Qr_1 \; \dots \; Qr_n] =QR.$$ 
\end{proof}
\Answer From Gram-Schmidt, we have an orthogonal basis $$\left\{\begin{bmatrix}1 \\ 1\\ 1\\ 1\\ \end{bmatrix}, \begin{bmatrix}-3\\ 1\\ 1\\ 1\\ \end{bmatrix}, \begin{bmatrix}0\\ -2\\ 1\\ 1\\ \end{bmatrix} \right\},$$
and scaling to get an orthonormal basis, we get $$\left\{\begin{bmatrix}1/2 \\ 1/2\\ 1/2\\ 1/2\\ \end{bmatrix}, \begin{bmatrix}-3/\sqrt{12}\\ 1/\sqrt{12} \\ 1/\sqrt{12}\\ 1/\sqrt{12} \\ \end{bmatrix}, \begin{bmatrix}0\\ -2\sqrt{6}\\ 1/\sqrt{6} \\ 1/\sqrt{6} \\ \end{bmatrix} \right\}.$$

Thus $Q=\begin{bmatrix}
1/2 & -3/\sqrt{12} & 0\\
1/2 & 1/\sqrt{12} & -2\/\sqrt{6}\\
1/2 & 1/\sqrt{12} & 1/\sqrt{6}\\
1/2  & 1/\sqrt{12} & 1/\sqrt{6}\\
\end{bmatrix}$. Since $A=QR$ and $Q$ has orthonormal columns, we get $R=Q^{T}A$. Thus $R=\begin{bmatrix}
2 & 3/2 & 1\\
0 & 3/\sqrt{12} & 2/\sqrt{12}\\
0 & 0 & 2/\sqrt{6}\\
\end{bmatrix}$.

\end{enumerate}  

\end{comment}

\end{document}


%%% Local ables: 
%%% mode: latex
%%% TeX-master: t
%%% End: 
